\documentclass[11pt,a4paper]{article}
\usepackage[margin=2.4cm]{geometry}
\usepackage{amsmath,amssymb,bm}
\usepackage{booktabs}
\usepackage{array}
\usepackage{framed}
\usepackage{xcolor}
\usepackage[colorlinks=true,linkcolor=blue,citecolor=blue]{hyperref}

\newcommand{\bO}{\bm{\Omega}}
\newcommand{\bw}{\bm{\omega}}
\newcommand{\bx}{\mathbf{x}}
\newcommand{\bu}{\mathbf{u}}
\newcommand{\ba}{\mathbf{a}}
\newcommand{\bb}{\mathbf{b}}
\newcommand{\bt}{\mathbf{t}}
\newcommand{\bj}{\mathbf{j}}
\newcommand{\bk}{\mathbf{k}}
\newcommand{\bq}{\mathbf{q}}
\newcommand{\bI}{\mathbf{I}}
\newcommand{\bi}{\mathbf{i}}
\newcommand{\It}{\tilde{\mathbf{I}}}
\newcommand{\bG}{\bm{\Gamma}}
\newcommand{\bn}{\hat{\mathbf{n}}}
\newcommand{\E}{\mathsf{E}}
\newcommand{\Id}{\mathsf{1}}
\newcommand{\R}{\mathsf{R}}
\newcommand{\Sm}{\mathsf{S}}
\newcommand{\K}{\mathsf{K}}
\newcommand{\dSL}{\Delta\mathrm{SL}}
\newcommand{\oc}{\mathcal{C}}
\newcommand{\dd}{\,\mathrm{d}}
\newcommand{\pM}{\partial M}

\definecolor{clarbar}{RGB}{70,110,170}
\renewenvironment{leftbar}{%
  \def\FrameCommand{{\color{clarbar}\vrule width 2.5pt}\hspace{8pt}}%
  \MakeFramed{\advance\hsize-\width\FrameRestore}}%
  {\endMakeFramed}
\newenvironment{clar}{\begin{leftbar}\small\noindent\textbf{Clarification.}\ }{\end{leftbar}}

\title{Rotational feedbacks in the elastic fingerprint problem:\\
the static theory of Al-Attar et al.\ (2024) and a dynamic extension}
\author{Working notes (prepared with Claude) --- to be checked}
\date{23 September 2026}

\begin{document}
\maketitle

\tableofcontents

\section{Purpose and conventions}

Section~\ref{sec:static} restates the theory of Al-Attar et al.\ (2024), hereafter AA24, needed for
rotational feedbacks. It covers the non-rotating loading problem and its reciprocity theorems (their
\S2, abridged where not needed) and the rotational theory in full (their \S3), together with the
spherically symmetric implementation (their \S4.1). The aim is that this section is a
\emph{superset} of AA24: every equation of theirs that is used appears with its original number,
and intermediate steps and remarks that they omit are added in shaded \textbf{Clarification}
boxes. Clarifications are my reading of the paper and should be checked.

Section~\ref{sec:dynamic} then replaces the static Euler equation with the linearised Liouville
(dynamic Euler) equation and repeats the derivation in the same style.

\paragraph{Equation tags.} Tags such as (3.12) reproduce AA24's numbering. Tags such as (3.12a)
are added intermediate steps. Tags such as (D3.15) are the dynamic counterparts of AA24's
equations. Tags such as (P1) concern the \texttt{pyslfp} implementation.

\paragraph{Notation.} This follows AA24 with one change: the ocean function is written $\oc$
rather than $C$, to avoid a clash with moments of inertia. Principal moments are
$I_1\le I_2<I_3$, as in AA24. The spherically symmetric implementation uses $I_1=I_2$. $\dSL$ is
the same quantity as $\xi$ in Heathcote et al.

\section{Static theory, following Al-Attar et al.\ (2024)}
\label{sec:static}

\subsection{Equations of motion for static loading (AA24 \S2.1)}

A non-rotating elastic earth model occupies volume $M$ with surface $\pM$, and is initially in
hydrostatic equilibrium with gravitational potential $\Phi$. Its response to a surface load
$\sigma$ is governed by the weak form (Al-Attar \& Tromp 2014, \S2.3.2)
\begin{equation}
  A(\bu,\phi\,|\,\bu',\phi')+\int_{\pM}(\bu'\cdot\nabla\Phi+\phi')\,\sigma\dd S=0 .
  \tag{2.1}\label{aa21}
\end{equation}
Here $\bu$ is the displacement, $\phi$ the Eulerian gravitational potential perturbation, and
$(\bu',\phi')$ test functions. $(\bu,\phi)$ is a solution if eq.~\eqref{aa21} holds for all test
functions. $A$ is the elastic--gravitational bilinear form (Al-Attar \& Tromp 2014, eq.~2.52) with
two properties:
\begin{equation}
  A(\bu,\phi\,|\,\bu',\phi')=A(\bu',\phi'\,|\,\bu,\phi),
  \tag{2.2}\label{aa22}
\end{equation}
where displacement and potential are swapped \emph{as pairs}, and
\begin{equation}
  A(\bu,\phi\,|\,\bu',\phi')=0 \quad\text{for all }(\bu',\phi')\quad\text{whenever}\quad
  \bu=\ba+\bb\times\bx,\ \ \phi=-(\ba+\bb\times\bx)\cdot\nabla\Phi,
  \tag{2.3--2.4}\label{aa24}
\end{equation}
for arbitrary constant $\ba,\bb$ (linearised rigid-body motions). Taking test functions of that
form, eq.~(2.5), eq.~\eqref{aa21} reduces to a trivial identity. By the Fredholm alternative the
problem is therefore solvable for any $\sigma$, with the solution unique up to a rigid-body motion.

\begin{clar}
The surface term is the virtual work of the load. A surface mass density $\sigma$ feels the
gravitational force $-\sigma\nabla\Phi$ per unit area, paired with $\bu'$. It is also a source in
Poisson's equation for $\phi$, paired with $\phi'$. With a rigid-body test function,
$\bu'\cdot\nabla\Phi+\phi'=\bu'\cdot\nabla\Phi-\bu'\cdot\nabla\Phi=0$ pointwise, so the surface term
vanishes for \emph{every} $\sigma$. This is why no compatibility condition on $\sigma$ arises. On
the surface of a spherically symmetric model, $\nabla\Phi=g\,\bn$, so
$\bu\cdot\nabla\Phi=g\,u_r$.
\end{clar}

\subsection{A first reciprocity theorem (AA24 \S2.2)}

Let $(\bu,\phi)$ solve the problem for $\sigma$, and $(\bu^\dagger,\phi^\dagger)$ for $\sigma^\dagger$.
Using each solution as the test function in the other problem gives
\begin{align}
  A(\bu,\phi\,|\,\bu^\dagger,\phi^\dagger)+\int_{\pM}(\bu^\dagger\cdot\nabla\Phi+\phi^\dagger)\,\sigma\dd S&=0,\tag{2.6}\\
  A(\bu^\dagger,\phi^\dagger\,|\,\bu,\phi)+\int_{\pM}(\bu\cdot\nabla\Phi+\phi)\,\sigma^\dagger\dd S&=0.\tag{2.7}
\end{align}
Subtracting and using eq.~\eqref{aa22},
\begin{equation}
  \int_{\pM}(\bu^\dagger\cdot\nabla\Phi+\phi^\dagger)\,\sigma\dd S
  =\int_{\pM}(\bu\cdot\nabla\Phi+\phi)\,\sigma^\dagger\dd S ,
  \tag{2.8}\label{aa28}
\end{equation}
which is the reciprocity theorem of Tromp \& Mitrovica (1999).

\subsection{Reformulation in terms of sea level (AA24 \S2.3)}

In hydrostatic sea-level theory,
\begin{equation}
  \dSL=-\frac{1}{g}(\bu\cdot\nabla\Phi+\phi)+\frac{\Phi_g}{g},
  \tag{2.9}\label{aa29}
\end{equation}
where the constant $\Phi_g$ is fixed by conservation of mass,
\begin{equation}
  \int_{\pM}\sigma\dd S=0,\qquad \sigma=\rho_w\oc\,\dSL+\zeta ,
  \tag{2.10--2.11}\label{aa211}
\end{equation}
with direct load $\zeta$. For ice, $\zeta=\rho_i(1-\oc)\Delta I$ (2.12). For two solutions sharing
the same ocean function,
\begin{equation}
  \bu\cdot\nabla\Phi+\phi=-g\,\dSL+\Phi_g,\qquad
  \bu^\dagger\cdot\nabla\Phi+\phi^\dagger=-g\,\dSL^\dagger+\Phi_g^\dagger,
  \tag{2.13}
\end{equation}
and eq.~\eqref{aa28} becomes
\begin{equation}
  \int_{\pM}(-g\,\dSL^\dagger+\Phi_g^\dagger)\,\sigma\dd S=\int_{\pM}(-g\,\dSL+\Phi_g)\,\sigma^\dagger\dd S .
  \tag{2.14}
\end{equation}
The constant terms vanish by mass conservation, giving
$\int_{\pM}\dSL^\dagger\sigma\dd S=\int_{\pM}\dSL\,\sigma^\dagger\dd S$ (2.15). Substituting the
load decompositions (2.16) and cancelling the symmetric ocean terms,
\begin{equation}
  \int_{\pM}\dSL^\dagger\,\zeta\dd S=\int_{\pM}\dSL\,\zeta^\dagger\dd S .
  \tag{2.17}\label{aa217}
\end{equation}

\begin{clar}
(i) $\int\Phi_g^\dagger\sigma\dd S=\Phi_g^\dagger\int\sigma\dd S=0$ because $\Phi_g^\dagger$ is a
constant. (ii) The cancelled terms are
$\rho_w\int\oc\,\dSL^\dagger\dSL\dd S$ on both sides. This requires the \emph{same} $\oc$ in both
problems, i.e.\ fixed shorelines. (iii) The factor $-g$ is common to both sides and divides out.
\end{clar}

\paragraph{Green's function and kernels (AA24 \S\S2.4--2.6).} Writing
$\dSL(\bx)=\int_{\pM}G(\bx,\bx')\zeta(\bx')\dd S_{\bx'}$ (2.18) and taking $\zeta^\dagger=\delta(\cdot,\bx)$
in eq.~\eqref{aa217} gives $G(\bx,\bx')=G(\bx',\bx)$ (2.21). For a functional $J$ with
\begin{equation}
  \delta J=\int_{\pM}\zeta^\dagger\,\delta\dSL\dd S,
  \tag{2.25}
\end{equation}
eq.~\eqref{aa217} gives
\begin{equation}
  \delta J=\int_{\pM}\dSL^\dagger\,\delta\zeta\dd S .
  \tag{2.26}
\end{equation}
Here $\dSL^\dagger$ (the adjoint sea level) solves the ordinary fingerprint problem with direct load
$\zeta^\dagger$ (the adjoint load). The ice-thickness kernel is $K=\rho_i(1-\oc)\dSL^\dagger$ (2.28).

\subsection{The generalised loading problem (AA24 \S2.7)}
\label{sec:gen_static}

For observables that depend directly on $\bu$ or $\phi$,
\begin{equation}
  A(\bu,\phi\,|\,\bu',\phi')+\int_{\pM}(\bu'\cdot\nabla\Phi+\phi')\,\sigma\dd S
  +\int_{\pM}(\bu'\cdot\bt+\zeta_\phi\,\phi')\dd S=0,
  \tag{2.29}\label{aa229}
\end{equation}
which is solvable if and only if
\begin{equation}
  \int_{\pM}(\bt-\zeta_\phi\nabla\Phi)\cdot(\ba+\bb\times\bx)\dd S=0
  \tag{2.30}\label{aa230}
\end{equation}
for all $\ba,\bb$. Pairing solutions for $(\sigma,\bt,\zeta_\phi)$ and
$(\sigma^\dagger,\bt^\dagger,\zeta^\dagger_\phi)$,
\begin{align}
  \int_{\pM}(\bu^\dagger\cdot\nabla\Phi+\phi^\dagger)\sigma\dd S+\int_{\pM}(\bt\cdot\bu^\dagger+\zeta_\phi\phi^\dagger)\dd S
  &=\int_{\pM}(\bu\cdot\nabla\Phi+\phi)\sigma^\dagger\dd S+\int_{\pM}(\bt^\dagger\cdot\bu+\zeta^\dagger_\phi\phi)\dd S,
  \tag{2.31}\label{aa231}\\
  \int_{\pM}\dSL^\dagger\zeta\dd S-\frac{1}{g}\int_{\pM}(\bt\cdot\bu^\dagger+\zeta_\phi\phi^\dagger)\dd S
  &=\int_{\pM}\dSL\,\zeta^\dagger\dd S-\frac{1}{g}\int_{\pM}(\bt^\dagger\cdot\bu+\zeta^\dagger_\phi\phi)\dd S.
  \tag{2.32}\label{aa232}
\end{align}
For a functional $J(\dSL,\bu,\phi)$, write its first-order perturbation as
\begin{equation}
  \delta J=\int_{\pM}\zeta^\dagger\,\delta\dSL\dd S-\frac1g\int_{\pM}(\bt^\dagger\cdot\delta\bu+\zeta^\dagger_\phi\,\delta\phi)\dd S
  \tag{2.33}\label{aa233}
\end{equation}
for suitable $(\zeta^\dagger,\bt^\dagger,\zeta^\dagger_\phi)$. Eq.~\eqref{aa232}, applied with the
forward problem $(\delta\dSL,\delta\zeta,\bt=\mathbf{0},\zeta_\phi=0)$, then gives
\begin{equation}
  \delta J=\int_{\pM}\dSL^\dagger\,\delta\zeta\dd S ,
  \tag{2.34}\label{aa234}
\end{equation}
where $\dSL^\dagger$ solves the generalised problem with loads $(\zeta^\dagger,\bt^\dagger,\zeta^\dagger_\phi)$.
It is well defined provided $J$ is invariant under rigid-body motions, AA24 eq.~(2.35).

\subsection{Review of rotational dynamics (AA24 \S3.1)}
\label{sec:aa31}

A model rotating steadily with angular velocity $\bO$ satisfies
\begin{equation}
  \bO\times(\bI\bO)=\mathbf{0},
  \tag{3.1}\label{aa31}
\end{equation}
so $\bO$ lies along a principal axis, taken to be that of the largest moment, $I_3$. After
deformation, $\bI\to\bI+\bi$ and $\bO\to\bO+\bw$. To first order, with $\bj=\bi\bO$,
\begin{equation}
  \bO\times\bj+I_3\,\bw\times\bO+\bO\times(\bI\bw)=\mathbf{0}.
  \tag{3.2}\label{aa32}
\end{equation}

\begin{clar}
Expanding $(\bO+\bw)\times[(\bI+\bi)(\bO+\bw)]$ to first order gives
$\bO\times(\bi\bO)+\bO\times(\bI\bw)+\bw\times(\bI\bO)$, and $\bI\bO=I_3\bO$. Taking
$\bO\times$ of eq.~\eqref{aa32} and using $\bO\times(\bO\times\mathbf{v})=\bO(\bO\cdot\mathbf{v})-\Omega^2\mathbf{v}=-\Omega^2\mathbf{v}_\perp$:
\begin{equation}
  -\Omega^2\bj_\perp+I_3\Omega^2\bw_\perp-\Omega^2(\bI\bw)_\perp=\mathbf{0}.
  \tag{3.2a}\label{aa32a}
\end{equation}
Here $\bO\times(\bw\times\bO)=\Omega^2\bw-\bO(\bO\cdot\bw)=\Omega^2\bw_\perp$ was used. Because
$\bO$ is along a principal axis, $\bI$ maps the equatorial plane to itself, so
$(\bI\bw)_\perp=\bI\bw_\perp$. This gives eq.~(3.3).
\end{clar}

\begin{equation}
  (I_3-\bI)\bw_\perp=\bj_\perp .
  \tag{3.3}\label{aa33}
\end{equation}
Equation~\eqref{aa32} contains no information about $\bw_\parallel$. Conservation of angular momentum,
\begin{equation}
  \|(\bI+\bi)(\bO+\bw)\|^2=\|\bI\bO\|^2,
  \tag{3.4}\label{aa34}
\end{equation}
gives to first order
\begin{equation}
  -I_3\,\bw_\parallel=\bj_\parallel .
  \tag{3.5}\label{aa35}
\end{equation}

\begin{clar}
(i) First order in eq.~\eqref{aa34}: $2(\bI\bO)\cdot(\bI\bw+\bi\bO)=0$. Since $\bI\bO=I_3\bO$ and
$\bI$ is symmetric, $(\bI\bO)\cdot(\bI\bw)=I_3^2\,\bO\cdot\bw$ and $(\bI\bO)\cdot\bj=I_3\bO\cdot\bj$.
Hence $I_3\bw_\parallel+\bj_\parallel=\mathbf{0}$. (ii) Only the \emph{magnitude} of the angular
momentum is used because $\mathbf{L}$ is fixed in inertial space, while the Earth-fixed frame in
which $\bw$ is measured may re-orient. $\|\mathbf{L}\|$ is invariant under that re-orientation.
Section~\ref{sec:dyn31} shows that eq.~\eqref{aa35} is also exactly the axial component of the
Liouville equation.
\end{clar}

Combining eqs.~\eqref{aa33} and \eqref{aa35},
\begin{equation}
  \It\bw=\bj,\qquad
  \It=\begin{pmatrix}I_3-I_1&0&0\\0&I_3-I_2&0\\0&0&-I_3\end{pmatrix}
  \tag{3.6--3.7}\label{aa36}
\end{equation}
in the principal-axis basis. $\It$ is symmetric and invertible. Because
$I_3\gg I_3-I_1\ge I_3-I_2$, the $x_3$ component of $\bw$ is commonly neglected; the reciprocity
results do not depend on this.

\begin{clar}
Invertibility needs $I_3>I_2$ strictly, i.e.\ rotation about the axis of greatest moment. The
symmetry of $\It$ is the property that makes reciprocity work in \S\ref{sec:aa32}. If $\bw_\parallel$
were retained, $\bw_\parallel=\omega_3\hat{\mathbf{x}}_3$ would give
$\psi=-\Omega\,\omega_3\,r^2\sin^2\theta$, which has degree-0 and degree-2, order-0 parts. The
equatorial components give degree~2, order~1 only (eq.~\ref{p1}).
\end{clar}

The centrifugal potential and its first-order perturbation are
\begin{equation}
  \Psi=-\tfrac12\|\bO\times\bx\|^2,\qquad \psi=-(\bO\times\bx)\cdot(\bw\times\bx),
  \tag{3.8--3.9}\label{aa39}
\end{equation}
with
\begin{equation}
  \nabla\Psi=\bO\times(\bO\times\bx),\qquad \nabla\psi=\bO\times(\bw\times\bx)+\bw\times(\bO\times\bx).
  \tag{3.10}\label{aa310}
\end{equation}

\begin{clar}
$\|\bO\times\bx\|^2=\Omega^2\|\bx\|^2-(\bO\cdot\bx)^2$, so
$\nabla\Psi=-\Omega^2\bx+(\bO\cdot\bx)\bO=\bO\times(\bO\times\bx)$. The centrifugal acceleration is
$-\nabla\Psi$. Replacing $\bO$ by $\bO+\bw$ and linearising gives $\psi$ and $\nabla\psi$.
\end{clar}

Next, the inertia perturbation must be related to the load. For an arbitrary vector $\bO$,
\begin{equation}
  \bO\cdot(\bI\bO)=\int_M\rho\,\|\bO\times\bx\|^2\dd V,
  \tag{3.11}\label{aa311}
\end{equation}
and under a surface load $\sigma$ with displacement $\bu$,
\begin{equation}
  \bO\cdot(\bi\bO)=-\int_M2\rho\,\bu\cdot[\bO\times(\bO\times\bx)]\dd V+\int_{\pM}\sigma\,\|\bO\times\bx\|^2\dd S .
  \tag{3.12}\label{aa312}
\end{equation}

\begin{clar}
(i) Eq.~\eqref{aa311} follows from $\bI=\int_M\rho(\|\bx\|^2\Id-\bx\bx^T)\dd V$. (ii) For
eq.~\eqref{aa312}, work in the reference configuration: mass is conserved and a particle moves from
$\bx$ to $\bx+\bu$, so the solid part contributes $\int_M\rho\|\bO\times(\bx+\bu)\|^2\dd V$. To first
order,
\begin{equation*}
  \|\bO\times(\bx+\bu)\|^2-\|\bO\times\bx\|^2\approx2(\bO\times\bx)\cdot(\bO\times\bu)=2\bu\cdot[(\bO\times\bx)\times\bO]=-2\bu\cdot[\bO\times(\bO\times\bx)].
\end{equation*}
The surface load adds mass $\sigma\dd S$ at $\bx\in\pM$. (iii) For a fluid outer core the Lagrangian
$\bu$ is not well defined; AA24 note that an Eulerian treatment there (Dahlen 1974) leaves the
results unchanged.
\end{clar}

Replacing $\bO$ by $\bO+\bw'$ in eq.~\eqref{aa312} and keeping the part linear in $\bw'$,
\begin{equation}
  \bw'\cdot\bj=-\int_M\rho\,\bu\cdot[\bw'\times(\bO\times\bx)+\bO\times(\bw'\times\bx)]\dd V
  +\int_{\pM}\sigma\,(\bO\times\bx)\cdot(\bw'\times\bx)\dd S,
  \tag{3.13}\label{aa313}
\end{equation}
and comparing with eqs.~\eqref{aa39}--\eqref{aa310},
\begin{equation}
  \boxed{\;\bw'\cdot\bj+\int_M\rho\,\bu\cdot\nabla\psi'\dd V+\int_{\pM}\sigma\,\psi'\dd S=0\;}
  \tag{3.14}\label{aa314}
\end{equation}
where $\psi'$ is the centrifugal perturbation associated with $\bw'$, while $\bj$ is associated
with $(\sigma,\bu)$.

\begin{clar}
In eq.~\eqref{aa312} the $\bO$ on the \emph{left} is a dummy vector, but $\bi$ does not depend on
it. Replacing $\bO\to\bO+\bw'$, the linear part of the left side is
$\bw'\cdot\bi\bO+\bO\cdot\bi\bw'=2\bw'\cdot\bj$, using the symmetry of $\bi$. The right side gives
twice the integrands shown, and the factor 2 cancels. Eq.~\eqref{aa314} is purely a statement about
the mass distribution at one instant. This is why it carries over unchanged to the dynamic problem
in \S\ref{sec:dynamic}.
\end{clar}

\begin{clar}
In a spherically symmetric calculation, $\bI$ is taken to be the \emph{observed} inertia tensor
(Mitrovica et al.\ 2005), even though the reference model is spherical. More generally, a rotating
hydrostatic reference state would have $\nabla(\Phi+\Psi)$ in place of $\nabla\Phi$ and
corresponding terms in $A$. As I read it, AA24 keep $\Psi$ only to generate $\psi$; the symmetry
arguments are the same either way.
\end{clar}

\subsection{Reciprocity theorems including rotational feedbacks (AA24 \S3.2)}
\label{sec:aa32}

The static loading problem in a rotating model, in weak form, is
\begin{equation}
  \boxed{\;A(\bu,\phi\,|\,\bu',\phi')+\int_{\pM}(\bu'\cdot\nabla\Phi+\phi')\,\sigma\dd S
  +\int_M\rho\,\bu'\cdot\nabla\psi\dd V+(\It\bw-\bj)\cdot\bw'=0\;}
  \tag{3.15}\label{aa315}
\end{equation}
for all $(\bu',\phi',\bw')$. The $\psi$ term is a tidal body force. The test function $\bw'$
enforces eq.~\eqref{aa36}; $\bj$ depends implicitly on $\bu$ and $\sigma$.

\begin{clar}
(i) The body force per unit volume is $-\rho\nabla\psi$. Its virtual work enters with the same
sign convention as the load term in eq.~\eqref{aa21}. (ii) The unknowns are now
$(\bu,\phi,\bw)$, and $\bw'$ acts as a Lagrange multiplier for the rotational constraint. (iii) The
problem is coupled: $\psi$ depends on $\bw$, $\bj$ depends on $\bu$ and $\sigma$, and $\sigma$
depends on $\dSL$, which depends on $\psi$ (eq.~3.18). This is why it is solved iteratively.
\end{clar}

Let $(\bu^\dagger,\phi^\dagger,\bw^\dagger)$ solve the problem for $\sigma^\dagger$. Then
\begin{equation}
  \int_{\pM}(\bu^\dagger\cdot\nabla\Phi+\phi^\dagger)\sigma\dd S+\int_M\rho\,\bu^\dagger\cdot\nabla\psi\dd V+\bj^\dagger\cdot\bw
  =\int_{\pM}(\bu\cdot\nabla\Phi+\phi)\sigma^\dagger\dd S+\int_M\rho\,\bu\cdot\nabla\psi^\dagger\dd V+\bj\cdot\bw^\dagger .
  \tag{3.16}\label{aa316}
\end{equation}

\begin{clar}
Write eq.~\eqref{aa315} for $\sigma$ with test functions $(\bu^\dagger,\phi^\dagger,\bw^\dagger)$,
and for $\sigma^\dagger$ with test functions $(\bu,\phi,\bw)$, then subtract. The $A$ terms cancel
by eq.~\eqref{aa22}, and $\bw^\dagger\cdot\It\bw=\bw\cdot\It\bw^\dagger$ cancels by symmetry of
$\It$. What remains is
$[\ldots]-\bj\cdot\bw^\dagger=[\ldots]^\dagger-\bj^\dagger\cdot\bw$. Moving the $\bj$ terms across
gives eq.~\eqref{aa316}.
\end{clar}

Using eq.~\eqref{aa314} twice,
\begin{equation}
  \boxed{\;\int_{\pM}(\bu^\dagger\cdot\nabla\Phi+\phi^\dagger+\psi^\dagger)\,\sigma\dd S=
  \int_{\pM}(\bu\cdot\nabla\Phi+\phi+\psi)\,\sigma^\dagger\dd S\;}
  \tag{3.17}\label{aa317}
\end{equation}

\begin{clar}
Apply eq.~\eqref{aa314} with $(\bj,\bu,\sigma)$ from the $\dagger$ problem and $\psi'=\psi$:
$\bj^\dagger\cdot\bw=-\int\rho\,\bu^\dagger\cdot\nabla\psi\dd V-\int\sigma^\dagger\psi\dd S$.
Similarly
$\bj\cdot\bw^\dagger=-\int\rho\,\bu\cdot\nabla\psi^\dagger\dd V-\int\sigma\psi^\dagger\dd S$.
Substituting into eq.~\eqref{aa316}, both volume integrals cancel on each side, and the surface
terms rearrange into eq.~\eqref{aa317}.
\end{clar}

With rotation, the sea surface is an equipotential of $\Phi+\Psi$, so
\begin{equation}
  \dSL=-\frac{1}{g}(\bu\cdot\nabla\Phi+\phi+\psi)+\frac{\Phi_g}{g},
  \tag{3.18}\label{aa318}
\end{equation}
and the argument of \S2.3 gives
\begin{equation}
  \boxed{\;\int_{\pM}\dSL^\dagger\,\zeta\dd S=\int_{\pM}\dSL\,\zeta^\dagger\dd S\;}
  \tag{3.19}\label{aa319}
\end{equation}
in the same form as eq.~\eqref{aa217}. Symmetry of the Green's function and the kernel results of
\S2.3 therefore carry over unchanged.

\subsection{The generalised problem with rotation (AA24 \S3.2, continued)}
\label{sec:gen_rot}

\subsubsection{What the generalised problem is, and why it is needed}

Eq.~\eqref{aa319} answers one question: how does \emph{sea level} respond to a \emph{direct load}?
Its practical use is through eq.~\eqref{aa234}: if a functional $J$ of the solution can be written
in the form of eq.~\eqref{aa233}, then its sensitivity kernel with respect to $\zeta$ is obtained by
solving one more fingerprint problem. The catch is that eq.~\eqref{aa233} must be expressible using
the available forcings. Most real observables are not functionals of $\dSL$ alone:
\begin{itemize}
  \item GRACE measures coefficients of the gravitational potential perturbation, $\phi$;
  \item GNSS measures vertical displacement, $\bn\cdot\bu$;
  \item satellite altimetry measures sea-surface height, which involves $\dSL$, $\bn\cdot\bu$ and (through the convention used to remove pole-tide effects) the centrifugal potential $\psi$, and hence $\bw$;
  \item polar-motion data constrain $\bw$ directly.
\end{itemize}
The generalised loading problem exists so that each of these can appear on the right-hand side of a
reciprocity theorem. The idea is simple: in the weak form, each unknown field is paired with its own
test function, so we add one forcing term conjugate to each unknown. Starting from
eq.~\eqref{aa315}, the three additions are
\begin{center}
\begin{tabular}{@{}llp{7.6cm}@{}}
\toprule
Forcing & Weak-form term & What it is\\
\midrule
$\bt$ & $\int_{\pM}\bt\cdot\bu'\dd S$ & a surface force density, conjugate to displacement\\
$\zeta_\phi$ & $\int_{\pM}\zeta_\phi\,\phi'\dd S$ & a fictitious surface mass: it sources the potential but exerts no force\\
$\bk$ & $\bk\cdot\bw'$ & a torque-like vector, conjugate to the rotation vector\\
\bottomrule
\end{tabular}
\end{center}
A genuine surface load $\sigma$ is the special case in which the first two act together in the fixed
combination $(\bu'\cdot\nabla\Phi+\phi')\sigma$: real mass both presses down on the surface and
sources the gravitational potential. Separating them is what allows an observable that sees only one
of the two --- for example a potential coefficient, which sees mass but feels no force --- to be
represented. AA24 note that $\bt$ and $\zeta_\phi$ ``need not have a direct physical
interpretation''; they are bookkeeping devices for the adjoint problem. In
\texttt{pyslfp} they are the arguments \texttt{displacement\_load}, \texttt{gravitational\_potential\_load}
and \texttt{angular\_momentum\_change} of the generalised solver, and in the spectral implementation
they act through the generalised Love numbers of eq.~(4.4).

\subsubsection{The equations}

The generalised static problem with rotation is
\begin{equation}
  A(\bu,\phi\,|\,\bu',\phi')+\int_{\pM}(\bu'\cdot\nabla\Phi+\phi')\sigma\dd S+\int_{\pM}(\bt\cdot\bu'+\zeta_\phi\phi')\dd S
  +\int_M\rho\,\bu'\cdot\nabla\psi\dd V+(\It\bw-\bj+\bk)\cdot\bw'=0 .
  \tag{3.20}\label{aa320}
\end{equation}
Because $\zeta_\phi$ acts as mass as far as the potential is concerned, it also contributes to the
inertia perturbation, and eq.~\eqref{aa314} generalises to
\begin{equation}
  \bw'\cdot\bj+\int_M\rho\,\bu\cdot\nabla\psi'\dd V+\int_{\pM}(\sigma+\zeta_\phi)\,\psi'\dd S=0 .
  \tag{3.21}\label{aa321}
\end{equation}

\begin{clar}
As I read it, $\zeta_\phi$ pairs with $\phi'$ exactly as $\sigma$ does, so it enters Poisson's
equation as a surface mass, but it is absent from the force term, so it does not load the solid
Earth. The inertia perturbation inferred from the perturbed potential therefore includes
$\int\zeta_\phi\|\bO\times\bx\|^2\dd S$. This is what the spectral implementation does
automatically, since $\bw_\perp$ is computed from $\phi_{21}$ (eq.~\ref{p2}).
\end{clar}

Repeating the argument of \S\ref{sec:aa32} gives the generalised reciprocity theorem
\begin{multline}
  \int_{\pM}\dSL^\dagger\zeta\dd S-\frac1g\int_{\pM}[\bt\cdot\bu^\dagger+\zeta_\phi(\phi^\dagger+\psi^\dagger)]\dd S-\frac1g\,\bk\cdot\bw^\dagger\\
  =\int_{\pM}\dSL\,\zeta^\dagger\dd S-\frac1g\int_{\pM}[\bt^\dagger\cdot\bu+\zeta^\dagger_\phi(\phi+\psi)]\dd S-\frac1g\,\bk^\dagger\cdot\bw .
  \tag{3.22}\label{aa322}
\end{multline}

\begin{clar}
Subtracting the two instances of eq.~\eqref{aa320} gives the analogue of eq.~\eqref{aa316}, with
extra terms $\int(\bt\cdot\bu^\dagger+\zeta_\phi\phi^\dagger)+\bk\cdot\bw^\dagger$ on the left and
their $\dagger$-counterparts on the right. Using eq.~\eqref{aa321} in place of eq.~\eqref{aa314}
produces the extra surface terms $\int\zeta_\phi\psi^\dagger$ and $\int\zeta_\phi^\dagger\psi$,
which is why $\zeta_\phi$ ends up paired with $\phi^\dagger+\psi^\dagger$ rather than
$\phi^\dagger$ alone. Then substitute eq.~\eqref{aa318}, use $\int\sigma\dd S=0$, divide by $-g$,
and cancel the symmetric ocean terms.
\end{clar}

Accordingly, if a functional $J(\dSL,\bu,\phi,\bw)$ has
\begin{equation}
  \delta J=\int_{\pM}\zeta^\dagger\,\delta\dSL\dd S-\frac1g\int_{\pM}[\bt^\dagger\cdot\delta\bu+\zeta^\dagger_\phi(\delta\phi+\delta\psi)]\dd S-\frac1g\,\bk^\dagger\cdot\delta\bw ,
  \tag{3.22a}\label{aa322a}
\end{equation}
then eq.~\eqref{aa322}, applied with a forward problem carrying only a direct load, gives
\begin{equation}
  \delta J=\int_{\pM}\dSL^\dagger\,\delta\zeta\dd S ,
  \tag{3.22b}\label{aa322b}
\end{equation}
where $\dSL^\dagger$ comes from solving eq.~\eqref{aa320} with the adjoint loads
$(\zeta^\dagger,\bt^\dagger,\zeta^\dagger_\phi,\bk^\dagger)$ read off from eq.~\eqref{aa322a}.

Finally, using $\psi=-[\bx\times(\bO\times\bx)]\cdot\bw$ (a scalar triple product applied to
eq.~\ref{aa39}), the $\zeta_\phi\psi$ terms can be moved into the $\bk$ terms:
\begin{multline}
  \int_{\pM}\dSL^\dagger\zeta\dd S-\frac1g\int_{\pM}(\bt\cdot\bu^\dagger+\zeta_\phi\phi^\dagger)\dd S-\frac1g\Big(\bk-\int_{\pM}\zeta_\phi\,\bx\times(\bO\times\bx)\dd S\Big)\cdot\bw^\dagger\\
  =\int_{\pM}\dSL\,\zeta^\dagger\dd S-\frac1g\int_{\pM}(\bt^\dagger\cdot\bu+\zeta^\dagger_\phi\phi)\dd S-\frac1g\Big(\bk^\dagger-\int_{\pM}\zeta^\dagger_\phi\,\bx\times(\bO\times\bx)\dd S\Big)\cdot\bw .
  \tag{3.23}\label{aa323}
\end{multline}
The point of this form is that $\zeta_\phi$ now multiplies $\phi$ alone. An observable in the bare
gravitational potential, such as a GRACE coefficient, is therefore represented using
eq.~\eqref{aa323} rather than eq.~\eqref{aa322}: one chooses $\zeta_\phi^\dagger$ from the
$\phi$-dependence and then uses $\bk^\dagger$ to cancel the unwanted $\psi$ contribution, as in the
third row of the table below.

\subsubsection{Worked examples: adjoint loads for common observables}

The following follow AA24 \S\S4.3--4.6. In each case the adjoint loads are read off by comparing
$\delta J$ with eq.~\eqref{aa322a}, or with the corresponding form implied by eq.~\eqref{aa323}.

\begin{center}
\begin{tabular}{@{}lllll@{}}
\toprule
Observable $J$ & $\zeta^\dagger$ & $\bt^\dagger$ & $\zeta_\phi^\dagger$ & $\bk^\dagger$\\
\midrule
$\dSL(\bx_0)$ & $\delta(\bx,\bx_0)$ & $\mathbf{0}$ & $0$ & $\mathbf{0}$\\[2pt]
$\bn\cdot\bu(\bx_0)$ & $0$ & $-g\,\delta(\bx,\bx_0)\,\bn$ & $0$ & $\mathbf{0}$\\[2pt]
$\phi_{lm}$ (potential coefficient) & $0$ & $\mathbf{0}$ & $-g\,Y_{lm}/b^2$ & $\int_{\pM}\zeta^\dagger_\phi\,\bx\times(\bO\times\bx)\dd S$\\[2pt]
$\int_{\pM}w\,\Delta\mathrm{SSH}\dd S$ & $w$ & $-g\,w\,\bn$ & $0$ & $\int_{\pM}w\,\bx\times(\bO\times\bx)\dd S$\\
\bottomrule
\end{tabular}
\end{center}

Note that $\psi$ enters $\dSL$ through eq.~\eqref{aa318}, so the first row needs no $\bk^\dagger$:
its dependence on $\bw$ is already carried by $\dSL$. A $\bk^\dagger$ is needed only when the
observable depends on $\bw$ \emph{separately} from sea level, as in the last two rows.

\paragraph{The altimetry row in detail.} Sea-surface height change is
$\Delta\mathrm{SSH}=\dSL+\bn\cdot\bu+\psi/g$ (AA24 eq.~4.21; the same as eq.~20 of Heathcote et al.),
where the last term mirrors the convention by which pole-tide effects are removed from altimetry
products. An estimate of global-mean sea level is formed by averaging with a weight $w$ supported
where data exist, so $J=\int_{\pM}w\,\Delta\mathrm{SSH}\dd S$ and
\begin{equation*}
  \delta J=\int_{\pM}w\,\delta\dSL\dd S+\int_{\pM}w\,\bn\cdot\delta\bu\dd S+\frac1g\int_{\pM}w\,\delta\psi\dd S .
\end{equation*}
Comparing term by term with eq.~\eqref{aa322a}: the first gives $\zeta^\dagger=w$; the second gives
$-\bt^\dagger/g=w\,\bn$, i.e.\ $\bt^\dagger=-g\,w\,\bn$; and for the third, writing
$\delta\psi=-[\bx\times(\bO\times\bx)]\cdot\delta\bw$,
\begin{equation*}
  \frac1g\int_{\pM}w\,\delta\psi\dd S=-\frac1g\Big[\int_{\pM}w\,\bx\times(\bO\times\bx)\dd S\Big]\cdot\delta\bw ,
\end{equation*}
which matches $-\bk^\dagger\cdot\delta\bw/g$ with $\bk^\dagger$ as tabulated. So the rotational
forcing $\bk$ exists precisely so that observables carrying a centrifugal correction, such as this
one, can be handled.

\begin{clar}
$\bx\times(\bO\times\bx)=\Omega\,[r^2\hat{\mathbf{x}}_3-\bx\,x_3]$, whose equatorial components are
$-\Omega x_1x_3$ and $-\Omega x_2x_3$ (degree 2, order 1) and whose axial component is
$\Omega(x_1^2+x_2^2)$ (degrees 0 and 2, order 0). Hence in the third row of the table $\bk^\dagger$
vanishes for every potential coefficient except $l=2$, $m=\pm1$ (and, if $\omega_3$ is retained,
$l=0$ and $l=2,m=0$). For all other coefficients, $\zeta^\dagger_\phi$ alone is required. The factor
$b^{-2}$ in $\zeta^\dagger_\phi$ assumes the coefficient is defined by
$\phi_{lm}=b^{-2}\int_{\pM}\phi\,Y_{lm}\dd S$ with orthonormal $Y_{lm}$; any other normalisation
simply rescales it.
\end{clar}

\begin{clar}
\emph{Solvability.} Eq.~\eqref{aa320} is solvable only if the generalised loads do no net work
against rigid-body motions, eq.~\eqref{aa230} in the non-rotating case. AA24 point out that this is
equivalent to requiring $J$ to be invariant under linearised rigid-body motions, i.e.\ to be a
physically meaningful observable. For instance $\bu(\bx_0)$ by itself is not (it changes under a
translation of the whole model), whereas $\bn\cdot\bu(\bx_0)$ relative to the centre of mass, sea
level, and potential coefficients with $l\ge2$ are.
\end{clar}

\subsection{Spherically symmetric implementation (AA24 \S4.1, plus \texttt{pyslfp})}
\label{sec:spectral}

With $\sigma=\sum_{lm}\sigma_{lm}Y_{lm}$ (4.1), using orthonormal real harmonics without the
Condon--Shortley phase,
\begin{equation}
  u_{lm}=h_l\sigma_{lm},\qquad\phi_{lm}=k_l\sigma_{lm},
  \tag{4.2}
\end{equation}
where $u$ is the vertical displacement. AA24's $k_l$ includes the potential of the load itself,
which is non-standard. For functionals depending on $\bu$ only through its vertical component,
$\bt=\zeta_u\nabla\Phi$ (4.3), and
\begin{equation}
  u_{lm}=h_l\sigma_{lm}+h^u_l\zeta^u_{lm}+h^\phi_l\zeta^\phi_{lm},\qquad
  \phi_{lm}=k_l\sigma_{lm}+k^u_l\zeta^u_{lm}+k^\phi_l\zeta^\phi_{lm},
  \tag{4.4}
\end{equation}
with $h_l=h_l^u+h_l^\phi$ and $k_l=k_l^u+k_l^\phi$ (4.5), and $g\,h^\phi_l=k^u_l$ (4.6).
Tidal Love numbers $h^T_l,k^T_l$ give the response to the body force $-\rho\nabla\psi$:
$u_{lm}=h^T_l\psi_{lm}$, $\phi_{lm}=k^T_l\psi_{lm}$, where $\phi$ excludes $\psi$ itself.

\begin{clar}
(4.5): comparing eq.~\eqref{aa21} with eq.~\eqref{aa229} and $\bt=\zeta_u\nabla\Phi$, a true load
$\sigma$ is equivalent to $\zeta_u=\zeta_\phi=\sigma$. (4.6): in eq.~\eqref{aa231} with
$\sigma=\sigma^\dagger=0$, take a unit $\zeta_u$ harmonic in one problem and a unit $\zeta_\phi$
harmonic in the other. Then $\int\bt\cdot\bu^\dagger=g\,h^\phi_l$ and
$\int\zeta^\dagger_\phi\phi=k^u_l$.
\end{clar}

\paragraph{Rotational terms (not written out in AA24).} Write $\bm{\phi}_{21}=(\phi^c_{21},\phi^s_{21})^T$
for the $\cos\varphi$ and $\sin\varphi$ coefficients, with fields evaluated at radius $b$. With
$\bO=\Omega\hat{\mathbf{x}}_3$ and $\bw_\perp=(\omega_1,\omega_2)$, eq.~\eqref{aa39} gives
$\psi=\Omega x_3(x_1\omega_1+x_2\omega_2)=\Omega r^2\sin\theta\cos\theta\,(\omega_1\cos\varphi+\omega_2\sin\varphi)$.
Since $Y^c_{21}=\sqrt{15/4\pi}\sin\theta\cos\theta\cos\varphi$,
\begin{equation}
  \bm{\psi}_{21}=r\,\bw_\perp,\qquad r=\sqrt{\tfrac{4\pi}{15}}\;\Omega b^2 .
  \tag{P1}\label{p1}
\end{equation}
From MacCullagh's formula, the exterior potential of a mass distribution has
$\phi^c_{21}=\sqrt{12\pi/5}\;G\,\mathrm{i}_{13}/b^3$, and similarly $\phi^s_{21}$ with
$\mathrm{i}_{23}$. With $\bj_\perp=\Omega(\mathrm{i}_{13},\mathrm{i}_{23})$ and $I_1=I_2$,
\begin{equation}
  \frac{\bj_\perp}{I_3-I_1}=\beta\,\bm{\phi}_{21},\qquad
  \beta=\sqrt{\tfrac{5}{12\pi}}\;\frac{\Omega b^3}{G(I_3-I_1)} ,
  \tag{P2}\label{p2}
\end{equation}
where $\bm\phi_{21}=k_2\bm\sigma_{21}+k_2^T\bm\psi_{21}$. The static relation eq.~\eqref{aa320},
restricted to equatorial components, is
\begin{equation}
  \bw_\perp=\bq,\qquad \bq:=\beta\bm\phi_{21}-\frac{\bk_\perp}{I_3-I_1},
  \tag{P3}\label{p3}
\end{equation}
with $\omega_3$ neglected. This is exactly the \texttt{pyslfp} update: $r$ and $\beta$ are
\texttt{rotation\_factor} and \texttt{inertia\_factor}, and $\bk$ is \texttt{angular\_momentum\_change}.
The feedback strength is
\begin{equation}
  \kappa:=\beta r k_2^T=\frac{k_2^T}{k_s},\qquad k_s=\frac{3G(I_3-I_1)}{\Omega^2b^5},
  \tag{P4}\label{p4}
\end{equation}
with $k_s=0.951$ for the \texttt{pyslfp} constants and $k_2^T=0.2986$ (PREM), so $\kappa=0.314$.

\begin{clar}
\emph{MacCullagh step.} With orthonormal harmonics, the exterior potential is
\begin{equation*}
  \phi(r)=-G\sum_{m}\frac{4\pi}{5}\,r^{-3}\,Y_{2m}(\hat{\bx})\int_M\rho\,r'^2\,Y_{2m}(\hat{\bx}')\dd V',
\end{equation*}
and $\int\rho\,r'^2Y^c_{21}\dd V=\sqrt{15/4\pi}\int\rho\,x_1x_3\dd V=-\sqrt{15/4\pi}\,\mathrm{i}_{13}$. Hence
\begin{equation*}
  \phi^c_{21}=\frac{G}{b^{3}}\,\frac{4\pi}{5}\sqrt{\frac{15}{4\pi}}\;\mathrm{i}_{13}=\sqrt{\frac{12\pi}{5}}\;\frac{G\,\mathrm{i}_{13}}{b^3}.
\end{equation*}
Because $\phi$ is the potential of \emph{all} perturbed mass (internal deformation, $\sigma$, and
$\zeta_\phi$ in the generalised problem), eq.~\eqref{p2} automatically contains both terms of
eq.~\eqref{aa312} and the $\zeta_\phi$ term of eq.~\eqref{aa321}.
\emph{Iteration.} \texttt{pyslfp} alternates: compute $\bm\psi_{21}$ from the current $\bw$;
compute $u$ and $\phi$ from $\sigma$ and $\psi$; update $\bw_\perp\leftarrow\bq$; form $\dSL$ from
eq.~\eqref{aa318}; update $\sigma$.
\end{clar}

\section{Dynamic theory}
\label{sec:dynamic}

\subsection{What changes, and what does not}

In AA24 the rotation vector is slaved to the instantaneous inertia perturbation: given $\bj$, the
relation $\It\bw=\bj$ fixes $\bw$ at once. The rotation axis is assumed to move immediately to
wherever the current mass distribution puts it. In reality the rotation axis has angular momentum,
so it cannot jump: when the figure axis moves, the rotation axis circles around the new equilibrium
position. That free circling is the Chandler wobble, with an observed period of about 433 days.
Quasi-static theory is the limit in which the forcing is slow compared with that period.

The single change made here is therefore
\begin{equation*}
  \underbrace{\It\bw=\bj}_{\text{AA24, eq.~(3.6)}}\qquad\longrightarrow\qquad
  \underbrace{\It\bw+\bG\dot\bw=\bj}_{\text{eq.~(D3.6) below}},
\end{equation*}
where $\bG$ is a ``gyroscopic'' matrix defined in eq.~\eqref{d37}. Everything else in
\S\ref{sec:static} survives, because the elastic Earth and the equilibrium ocean respond
instantaneously: the only memory in the whole problem is in $\bw$. In detail:

\begin{center}
\begin{tabular}{@{}lll@{}}
\toprule
 & Static (AA24) & Dynamic (here)\\
\midrule
Rotation relation & $\It\bw=\bj$ & $\It\bw+\bG\dot\bw=\bj$\\
Axial component & $-I_3\omega_3=j_3$ & unchanged\\
Centrifugal potential & $\psi=-(\bO\times\bx)\cdot(\bw\times\bx)$ & unchanged, evaluated at each instant\\
Inertia identity (3.14) & holds & holds at each instant\\
Elastic and ocean response & instantaneous & unchanged\\
Free motion of $\bw$ & none & Chandler wobble\\
Reciprocity & $\int\dSL^\dagger\zeta=\int\dSL\,\zeta^\dagger$ & same, integrated over time\\
Adjoint problem & same equations & same equations, run backwards in time\\
Numerics & fixed-point iteration & same iteration, plus a 2-vector stepped in time\\
\bottomrule
\end{tabular}
\end{center}

Sections~\ref{sec:dyn31}--\ref{sec:dyn_gen} do the theory. Sections~\ref{sec:dyn_spec}--\ref{sec:dyn_num}
are the practical part and can be read on their own: eqs.~\eqref{dp1}--\eqref{dp2} are the whole
physical content, the table in \S\ref{sec:dyn_resp} says when it matters, and eq.~\eqref{dp8} with
the box in \S\ref{sec:dyn_num} is what a solver has to do differently.

\subsection{Rotational dynamics with time dependence}
\label{sec:dyn31}

In the Earth-fixed frame, with no external torque, the angular momentum
$\mathbf{L}=(\bI+\bi)(\bO+\bw)$ satisfies
\begin{equation}
  \frac{\dd}{\dd t}\Big[(\bI+\bi)(\bO+\bw)\Big]+(\bO+\bw)\times\Big[(\bI+\bi)(\bO+\bw)\Big]=\mathbf{0},
  \tag{D3.1}\label{d31}
\end{equation}
of which eq.~\eqref{aa31} is the steady, unperturbed case. To first order, with $\bj=\bi\bO$ and
$\bO$ constant,
\begin{equation}
  \boxed{\;\bI\dot\bw+\dot\bj+\bO\times\bj+I_3\,\bw\times\bO+\bO\times(\bI\bw)=\mathbf{0}\;}
  \tag{D3.2}\label{d32}
\end{equation}
which is eq.~\eqref{aa32} plus the two time-derivative terms.

\begin{clar}
(i) The time derivative is taken in the Earth-fixed frame, which is why the transport term
$(\bO+\bw)\times\mathbf{L}$ appears. (ii) Torque-free: surface mass redistribution is internal to
the Earth--ocean--ice system, and lunisolar torques act at nearly diurnal frequencies in this frame,
outside the band considered. (iii) The relative angular momentum of the deformation,
$\int\rho\,\bx\times\dot\bu\dd V$, is omitted; it vanishes in a Tisserand frame and is in any case
$O(\epsilon)$ (see below). Relative angular momentum of ocean currents or the atmosphere would enter
as $-\bk$ in the generalised problem.
\end{clar}

\paragraph{Equatorial components.} Taking $\bO\times$ of eq.~\eqref{d32} and reusing the steps of
eq.~\eqref{aa32a},
\begin{equation}
  \bO\times(\bI\dot\bw)+\bO\times\dot\bj+\Omega^2\Big[(I_3-\bI)\bw_\perp-\bj_\perp\Big]=\mathbf{0}.
  \tag{D3.3a}\label{d33a}
\end{equation}
Define the \emph{gyroscopic operator}
\begin{equation}
  \bG\mathbf{v}:=\frac{1}{\Omega^2}\,\bO\times(\bI\mathbf{v}),\qquad
  \bG=\frac{1}{\Omega}\begin{pmatrix}0&-I_2&0\\ I_1&0&0\\0&0&0\end{pmatrix},
  \tag{D3.3b}\label{d37}
\end{equation}
in the principal-axis basis. $\bG$ maps into the equatorial plane and annihilates axial vectors.
Then eq.~\eqref{d33a} is
\begin{equation}
  (I_3-\bI)\bw_\perp+\bG\dot\bw+\frac{1}{\Omega^2}\bO\times\dot\bj=\bj_\perp .
  \tag{D3.3}\label{d33}
\end{equation}

\paragraph{Ordering.} Suppose the fields vary at angular frequency $\nu$, that is with period
$2\pi/\nu$. Comparing the size of each term in eq.~\eqref{d33} with the static balance
$(I_3-\bI)\bw_\perp\sim\bj_\perp$:
\begin{itemize}
  \item $\bG\dot\bw\big/(I_3-I_1)\bw_\perp\sim I_1\nu/[\Omega(I_3-I_1)]=\nu/\sigma_e$, where
  $\sigma_e:=\Omega(I_3-I_1)/I_1$ is the Euler frequency, $2\pi/\sigma_e=304.5$\,d: \textbf{retain}.
  \item $\Omega^{-2}\bO\times\dot\bj\big/\bj_\perp\sim\nu/\Omega$: drop, for periods much longer than a day.
\end{itemize}
The two ratios differ by the factor $\Omega/\sigma_e=I_1/(I_3-I_1)\approx305$, which is why the
inertial term matters while everything else of its apparent order does not. Numerically,
$\nu/\sigma_e$ is $10.2$ at a period of one month, $0.83$ at one year and $0.08$ at ten years: the
dynamic term dominates the monthly band, is comparable to the static term at annual periods, and is
negligible for decadal trends.

\begin{clar}
$\epsilon:=1/(\Omega T)$ and $\delta:=(I_3-I_1)/I_1\approx3.3\times10^{-3}$ give the same statement:
the retained term is $O(\epsilon/\delta)$ and the dropped one $O(\epsilon)$. The retained term is
the only one amplified by the small dynamical ellipticity $\delta$.
\end{clar}

\paragraph{Axial component.} Dotting eq.~\eqref{d32} with $\bO$ removes every cross product. Using
$\bO\cdot\bI\dot\bw=(\bI\bO)\cdot\dot\bw=I_3\,\bO\cdot\dot\bw$,
\begin{equation}
  I_3\,\dot\bw_\parallel+\dot\bj_\parallel=\mathbf{0}
  \quad\Longrightarrow\quad
  -I_3\,\bw_\parallel=\bj_\parallel ,
  \tag{D3.5}\label{d35}
\end{equation}
given an unperturbed initial state. This is eq.~\eqref{aa35} exactly: the angular-momentum argument
of the static theory is simply the axial Liouville equation, and involves no approximation.

\begin{clar}
In eq.~\eqref{d35} the $\dot\bj_\parallel$ term is \emph{not} dropped: it is the only term balancing
$I_3\dot\bw_\parallel$, and the ordering argument applies to the equatorial balance only.
\end{clar}

Combining eqs.~\eqref{d33} and \eqref{d35}, and noting that $\bG$ has no axial part:
\begin{equation}
  \boxed{\;\It\bw+\bG\dot\bw=\bj\;}
  \tag{D3.6}\label{d36}
\end{equation}
This is the dynamic counterpart of eq.~\eqref{aa36}, which is recovered when $\dot\bw=\mathbf{0}$.
The key structural property is
\begin{equation}
  \bG^T=-\bG\quad\Longleftrightarrow\quad I_1=I_2 ,
  \tag{D3.7}\label{d37a}
\end{equation}
since $\bG=\Omega^{-2}[\bO]_\times\bI$ and $\bG^T=-\Omega^{-2}\bI[\bO]_\times$. From here on,
$I_1=I_2$. The antisymmetry of $\bG$ plays the same role in \S\ref{sec:dyn_recip} that the symmetry
of $\It$ played in \S\ref{sec:aa32}.

\paragraph{Centrifugal potential and inertia identity.} Eqs.~\eqref{aa39}--\eqref{aa310} hold at
each instant with $\bw=\bw(t)$. The Euler force $-\rho\,\dot\bw\times\bx$ in the momentum equation
is non-conservative, but it is $O(\epsilon)$ relative to $\rho\nabla\psi$ and is dropped; this is
what allows the elastic problem to remain forced by the potential $\psi$ alone, so that Love-number
methods still apply. Coriolis and inertial terms in the solid Earth, and the non-equilibrium ocean
response, are likewise $O(\epsilon)$. Eqs.~\eqref{aa311}--\eqref{aa314} involve only the mass
distribution at one instant, so
\begin{equation}
  \bw'\cdot\bj(t)+\int_M\rho\,\bu(t)\cdot\nabla\psi'\dd V+\int_{\pM}\sigma(t)\,\psi'\dd S=0
  \qquad\text{for each }t .
  \tag{D3.14}\label{d314}
\end{equation}

\subsection[Retaining the inertia-rate term]{Retaining the $\dot\bj$ term}
\label{sec:dyn_jdot}

The ordering argument above drops $\Omega^{-2}\bO\times\dot\bj$. It can instead be kept exactly. The
reason is that for an elastic Earth with an equilibrium ocean, $\bj$ is an \emph{instantaneous}
linear functional of the current load and the current rotation, with no memory: in the spectral form
of \S\ref{sec:spectral},
\begin{equation}
  \bj_\perp=A\,\mathbf{s}+B\,\bw_\perp,\qquad A=(I_3-I_1)\beta k_2,\qquad B=(I_3-I_1)\kappa,
  \tag{D3.8}\label{d38}
\end{equation}
where $\mathbf{s}=\bm\sigma_{21}$ and $\kappa=k_2^T/k_s$ as in eq.~\eqref{p4}. Hence
$\dot\bj_\perp=A\dot{\mathbf{s}}+B\dot\bw_\perp$, and eq.~\eqref{d33} can be written in full as
\begin{equation}
  (I_3-I_1-B)\,\bw_\perp+\frac{I_1+B}{\Omega}\,\E\dot\bw_\perp
  =A\Big(\mathbf{s}-\frac{1}{\Omega}\E\dot{\mathbf{s}}\Big),
  \tag{D3.9}\label{d39}
\end{equation}
with $\E$ the equatorial form of $\bG$ introduced in \S\ref{sec:dyn_spec}. Three things change
relative to eq.~\eqref{d36}: the restoring term is softened to $I_3-I_1-B$, which is the familiar
$1-\kappa$ factor; the effective inertia is stiffened to $I_1+B$, because the deformation adds to
the inertia of the wobble; and the forcing gains $-\E\dot{\mathbf{s}}/\Omega$, which is the $\dot c/\Omega$
term of the classical excitation function (Munk \& MacDonald 1960). The free frequency becomes
$\sigma_{\rm CW}=\Omega(I_3-I_1-B)/(I_1+B)$.

\paragraph{Angular momentum as the state variable.} Eq.~\eqref{d39} requires $\dot{\mathbf{s}}$.
Stepping the equatorial angular momentum perturbation instead avoids it: with
$\mathbf{l}=I_1\bw_\perp+\bj_\perp=(I_1+B)\bw_\perp+A\mathbf{s}$, the torque-free equation
$\dot{\mathbf{l}}=\Omega\E[(I_3-I_1)\bw_\perp-\bj_\perp]$ closes as
\begin{equation}
  \dot{\mathbf{l}}=\E\Big[\sigma_{\rm CW}\,\mathbf{l}-\frac{\Omega I_3 A}{I_1+B}\,\mathbf{s}\Big],
  \qquad \bw_\perp=\frac{\mathbf{l}-A\mathbf{s}}{I_1+B}.
  \tag{D3.10}\label{d310}
\end{equation}

\begin{clar}
Sizes: $B/I_1=\delta\kappa\approx1.0\times10^{-3}$, so the stiffening lengthens $T_{\rm CW}$ by about
half a day, well inside the uncertainty from the core and anelasticity. The $\dot{\mathbf{s}}$ term
is of relative size $\nu/\Omega$, that is (1 day)/(forcing period): 3\% at monthly periods, 0.3\% at
annual. The one qualitative effect is under a step load, where conservation of angular momentum
forces $\bw_\perp$ to jump instantly to $-A\mathbf{s}/(I_1+B)$ --- the rotation axis holding its
place in space while the figure axis moves beneath it --- before it begins to circle. That jump is
about 0.2\% of the eventual static response and of opposite sign.

$A$ and $B$ are pure Love-number expressions only for the solid Earth; once the ocean pole tide is
included, $B$ becomes the matrix $(I_3-I_1)\K$ of \S\ref{sec:dyn_spec}, which does not commute with
$\E$, and the stiffening is anisotropic. The construction also relies on the inertia having no
memory, so it does not carry over to a viscoelastic Earth, where $\bj$ is a convolution over the
load history.
\end{clar}

\paragraph{Reciprocity with $\dot\bj$ retained.} The reciprocity argument of
\S\ref{sec:dyn_recip} survives, and the quickest way to see why is to work with the operator
\begin{equation}
  \mathsf{X}:=\E\,\frac{\dd}{\dd t}
  \tag{D3.11}\label{d311}
\end{equation}
acting on time-dependent equatorial vectors under $\langle a,b\rangle=\int_0^T a\cdot b\dd t$.
Because $\E$ and $\dd/\dd t$ are each antisymmetric and act on different indices, $\mathsf{X}$ is
\emph{symmetric} up to boundary terms:
\begin{equation*}
  \langle \mathsf{X}a,b\rangle=-\int_0^T\dot a\cdot(\E b)\dd t
  =-\big[a\cdot\E b\big]_0^T+\langle a,\mathsf{X}b\rangle .
\end{equation*}
This is what eq.~\eqref{d317a} is really using: $\bG\,\dd/\dd t=(I_1/\Omega)\mathsf{X}$ is symmetric,
so the mismatch between the two problems integrates to a boundary term. The test for any additional
term is therefore just whether it is a multiple of $\mathsf{X}$.

Write $\mathbf{s}=\mathbf{s}^\zeta+\mathsf{S}\bw_\perp$, with $\mathsf{S}$ the ocean's degree-2,
order-1 load response to a unit rotation, and collect eq.~\eqref{d39} onto the left. The coefficient
of $\dot\bw_\perp$ is then
\begin{equation}
  \mathsf{G}=\E\big[(I_1+B)\Id+A\,\mathsf{S}\big],
  \tag{D3.12}\label{d312}
\end{equation}
and $\mathsf{G}\,\dd/\dd t$ is symmetric if and only if $\mathsf{G}$ is antisymmetric, that is iff
$[\mathsf{S},\E]=\mathbf{0}$; for a symmetric $2\times2$ matrix that means $\mathsf{S}\propto\Id$.
Equivalently, the forward operator $\mathsf{L}=\alpha\Id+\beta\mathsf{X}$ and the forcing operator
$\mathsf{F}=A(\Id-\mathsf{X}/\Omega)$ of eq.~\eqref{d39} are then both polynomials in $\mathsf{X}$
with scalar coefficients, so both are self-adjoint and they commute, making
$\mathsf{L}^{-1}\mathsf{F}$ symmetric.

\begin{clar}
So with no ocean, or an isotropic ocean response, eq.~\eqref{d319} holds unchanged when $\dot\bj$ is
retained, with the same terminal condition and the same time reversal. With a realistic anisotropic
$\mathsf{S}$ a residual asymmetry remains, of relative size
$|A\mathsf{S}|/(I_1+B)\approx\delta\,|\K_{\rm ocean}|\approx2\times10^{-4}$: two orders below the
$\nu/\Omega$ terms retained in eq.~\eqref{d39}, and far below the truncation of the theory itself.

This also tightens the triaxiality caveat into a single criterion covering every case:
\emph{every matrix multiplying $\dd/\dd t$ must be a multiple of $\E$}. That one condition accounts
for $I_1=I_2$, for the $B$ renormalisation, and for the ocean term.
\end{clar}

Everything that follows uses eq.~\eqref{d36}, i.e.\ with $\dot\bj$ dropped.

\subsection{Reciprocity with dynamic rotation}
\label{sec:dyn_recip}

\paragraph{What to expect.} The static problem is self-adjoint: the adjoint problem is the same
fingerprint problem driven by a different load. The dynamic problem has memory, so this cannot be
quite true --- a load applied today cannot affect sea level yesterday, and the sensitivity of a
measurement made today must run backwards over the loading history that preceded it. The result
below is that reciprocity survives in time-integrated form, with the adjoint problem being the same
equations solved \emph{backwards in time}.

On $t\in[0,T]$, the equations of motion in weak form are
\begin{equation}
  \boxed{\;A(\bu,\phi\,|\,\bu',\phi')+\int_{\pM}(\bu'\cdot\nabla\Phi+\phi')\,\sigma\dd S
  +\int_M\rho\,\bu'\cdot\nabla\psi\dd V+(\It\bw+\bG\dot\bw-\bj)\cdot\bw'=0,\qquad \bw(0)=\mathbf{0}\;}
  \tag{D3.15}\label{d315}
\end{equation}
for all $t$ and all test functions. At each instant, $(\bu,\phi)$ are determined by $\sigma(t)$ and
$\psi(t)$; the only memory is in $\bw$.

Let $(\bu^\dagger,\phi^\dagger,\bw^\dagger)$ solve eq.~\eqref{d315} for a load $\sigma^\dagger(t)$,
but with the \emph{terminal} condition $\bw^\dagger(T)=\mathbf{0}$ instead of an initial one. Then
for each $t$,
\begin{multline}
  \int_{\pM}(\bu^\dagger\cdot\nabla\Phi+\phi^\dagger)\sigma\dd S+\int_M\rho\,\bu^\dagger\cdot\nabla\psi\dd V+\bj^\dagger\cdot\bw+\bw^\dagger\cdot\bG\dot\bw\\
  =\int_{\pM}(\bu\cdot\nabla\Phi+\phi)\sigma^\dagger\dd S+\int_M\rho\,\bu\cdot\nabla\psi^\dagger\dd V+\bj\cdot\bw^\dagger+\bw\cdot\bG\dot\bw^\dagger ,
  \tag{D3.16}\label{d316}
\end{multline}
which is eq.~\eqref{aa316} plus the two $\bG$ terms. Using eq.~\eqref{d314} twice, exactly as for
eq.~\eqref{aa317},
\begin{equation}
  \int_{\pM}(\bu^\dagger\cdot\nabla\Phi+\phi^\dagger+\psi^\dagger)\sigma\dd S-\int_{\pM}(\bu\cdot\nabla\Phi+\phi+\psi)\sigma^\dagger\dd S
  =-\frac{\dd}{\dd t}\big(\bw^\dagger\cdot\bG\bw\big).
  \tag{D3.17a}\label{d317a}
\end{equation}

\begin{clar}
The left-hand side is what AA24 prove to be zero in the static case. The right-hand side comes from
the two $\bG$ terms: by eq.~\eqref{d37a},
$\bw\cdot\bG\dot\bw^\dagger=(\bG^T\bw)\cdot\dot\bw^\dagger=-\dot\bw^\dagger\cdot\bG\bw$, so their
difference is $-(\dot\bw^\dagger\cdot\bG\bw+\bw^\dagger\cdot\bG\dot\bw)$, a total time derivative.
This is the whole reason the antisymmetry of $\bG$ matters: it converts the mismatch into something
that integrates to a boundary term. For $I_1\neq I_2$ the symmetric part of $\bG$ leaves a
remainder that is not a total derivative, and the argument fails.
\end{clar}

Integrating over $[0,T]$, the boundary term $[\bw^\dagger\cdot\bG\bw]_0^T$ vanishes because
$\bw(0)=\mathbf{0}$ kills it at $t=0$ and $\bw^\dagger(T)=\mathbf{0}$ kills it at $t=T$:
\begin{equation}
  \boxed{\;\int_0^T\!\!\int_{\pM}(\bu^\dagger\cdot\nabla\Phi+\phi^\dagger+\psi^\dagger)\,\sigma\dd S\dd t=\int_0^T\!\!\int_{\pM}(\bu\cdot\nabla\Phi+\phi+\psi)\,\sigma^\dagger\dd S\dd t\;}
  \tag{D3.17}\label{d317}
\end{equation}
Eq.~\eqref{aa318} holds at each instant, and the steps of \S2.3 are instantaneous (the ocean terms
cancel pointwise in $t$), so
\begin{equation}
  \boxed{\;\int_0^T\!\!\int_{\pM}\dSL^\dagger\,\zeta\dd S\dd t=\int_0^T\!\!\int_{\pM}\dSL\,\zeta^\dagger\dd S\dd t\;}
  \tag{D3.19}\label{d319}
\end{equation}
This is AA24's eq.~\eqref{aa319} with an added time integral. The terminal condition
$\bw^\dagger(T)=\mathbf{0}$ is not an extra assumption we are free to drop: it is what the adjoint
problem must satisfy for the boundary term to vanish.

\paragraph{The adjoint problem, in practice.} The $\dagger$ fields satisfy the \emph{same}
equations as the forward problem, but with a condition imposed at the end rather than the start, so
they are integrated backwards from $T$. Writing $\tau=T-t$ and $\tilde\bw(\tau)=\bw^\dagger(T-\tau)$,
so that $\dd\tilde\bw/\dd\tau=-\dot\bw^\dagger$, eq.~\eqref{d36} becomes
\begin{equation}
  \It\tilde\bw-\bG\,\frac{\dd\tilde\bw}{\dd\tau}=\tilde\bj,\qquad\tilde\bw(0)=\mathbf{0}.
  \tag{D3.19a}\label{d319a}
\end{equation}
Here a tilde denotes an adjoint field re-expressed in $\tau$: $\tilde\bj(\tau)=\bj^\dagger(T-\tau)$
is the inertia perturbation of the adjoint problem read backwards in time, and likewise
$\tilde\zeta(\tau)=\zeta^\dagger(T-\tau)$ for the adjoint loads. In words: to compute an adjoint
field, take the adjoint loads, reverse them in time, and run the ordinary forward solver with the
sign of $\bG$ flipped.

\begin{clar}
$\bG\to-\bG$ is equivalent to reversing the sense of the reference rotation. Under $\bO\to-\bO$
together with $\bw\to-\bw$: $\psi$ is unchanged, $\It\bw\to-\It\bw$, $\bj=\bi\bO\to-\bj$, and
$\bG\dot\bw\to+\bG\dot\bw$ (both $\bG$ and $\dot\bw$ flip). Dividing by $-1$ gives
eq.~\eqref{d319a}, while $\dSL$ is unchanged. This is the familiar reciprocity of rotating systems:
source and receiver are exchanged and the rotation is reversed.
\end{clar}

\paragraph{Green's function.} By linearity and time invariance,
$\dSL(\bx,t)=\int_0^t\!\int_{\pM}G_{\bG}(\bx,\bx';t-t')\,\zeta(\bx',t')\dd S_{\bx'}\dd t'$, the
dynamic version of eq.~(2.18). Taking $\zeta=\delta(\cdot,\bx_1)\delta(t-t_1)$ and
$\zeta^\dagger=\delta(\cdot,\bx_2)\delta(t-t_2)$, and using eq.~\eqref{d319a} to identify
$\dSL^\dagger(\bx,t)=G_{-\bG}(\bx,\bx_2;t_2-t)$, eq.~\eqref{d319} gives
\begin{equation}
  G_{\bG}(\bx_2,\bx_1;s)=G_{-\bG}(\bx_1,\bx_2;s).
  \tag{D2.21}\label{d221}
\end{equation}
Swapping source and receiver is compensated by reversing the rotation. Integrating over $s$
recovers the static symmetry (2.21), because the time-integrated response does not depend on
$\bG$.

\paragraph{Sensitivity kernels.} If $\delta J=\int_0^T\!\int_{\pM}\zeta^\dagger\,\delta\dSL\dd S\dd t$, then
\begin{equation}
  \delta J=\int_0^T\!\!\int_{\pM}\dSL^\dagger\,\delta\zeta\dd S\dd t,\qquad
  K(\bx,t)=\rho_i(1-\oc)\,\dSL^\dagger(\bx,t),
  \tag{D2.26}\label{d226}
\end{equation}
so the kernel is now a function of position \emph{and} time: it says how much a load applied at
$\bx$ at time $t$ would have changed the observable. For an observation made at a single time
$t_o$, $\zeta^\dagger\propto\delta(t-t_o)$, and the backward solution vanishes for $t>t_o$, so the
kernel is automatically zero for loads applied after the measurement.

\paragraph{Non-zero initial state.} If $\bw(0)=\bw_0$, representing a free wobble already present
when the calculation starts (with $\bw^\dagger(T)=\mathbf{0}$ still), the boundary term at $t=0$
survives and eq.~\eqref{d319} becomes
\begin{equation}
  \int_0^T\!\!\int_{\pM}\dSL\,\zeta^\dagger\dd S\dd t=\int_0^T\!\!\int_{\pM}\dSL^\dagger\,\zeta\dd S\dd t+\frac1g\,\bw^\dagger(0)\cdot\bG\bw_0 ,
  \tag{D3.19b}\label{d319b}
\end{equation}
so $\partial J/\partial\bw_0=g^{-1}\bG^T\bw^\dagger(0)$. In an inversion, this is the kernel one
needs if the initial free wobble is treated as an unknown nuisance parameter alongside the loads.

\begin{clar}
From eq.~\eqref{d317a}, the space--time difference equals
$-[\bw^\dagger\cdot\bG\bw]_0^T=\bw^\dagger(0)\cdot\bG\bw_0$. Using eq.~\eqref{aa318} and mass
conservation, the left side equals
$-g\big[\int\!\!\int\dSL^\dagger\zeta-\int\!\!\int\dSL\,\zeta^\dagger\big]$. Divide by $-g$ and
rearrange.
\end{clar}

\subsection{The generalised dynamic problem}
\label{sec:dyn_gen}

Everything in \S\ref{sec:gen_rot} carries over with the loads now functions of time. The problem is
\begin{equation}
  A(\bu,\phi\,|\,\bu',\phi')+\int_{\pM}(\bu'\cdot\nabla\Phi+\phi')\sigma\dd S+\int_{\pM}(\bt\cdot\bu'+\zeta_\phi\phi')\dd S
  +\int_M\rho\,\bu'\cdot\nabla\psi\dd V+(\It\bw+\bG\dot\bw-\bj+\bk)\cdot\bw'=0,
  \tag{D3.20}\label{d320}
\end{equation}
with eq.~\eqref{aa321} holding pointwise in $t$, and the reciprocity theorem is eq.~\eqref{aa322}
integrated over $[0,T]$:
\begin{multline}
  \int_0^T\!\Big\{\int_{\pM}\dSL^\dagger\zeta\dd S-\frac1g\int_{\pM}[\bt\cdot\bu^\dagger+\zeta_\phi(\phi^\dagger+\psi^\dagger)]\dd S-\frac1g\,\bk\cdot\bw^\dagger\Big\}\dd t\\
  =\int_0^T\!\Big\{\int_{\pM}\dSL\,\zeta^\dagger\dd S-\frac1g\int_{\pM}[\bt^\dagger\cdot\bu+\zeta^\dagger_\phi(\phi+\psi)]\dd S-\frac1g\,\bk^\dagger\cdot\bw\Big\}\dd t ,
  \tag{D3.22}\label{d322}
\end{multline}
for $\bw(0)=\mathbf{0}$ and $\bw^\dagger(T)=\mathbf{0}$; the analogue of eq.~\eqref{aa323} follows
in the same way. The adjoint loads are exactly those tabulated in \S\ref{sec:gen_rot}, now carrying
whatever time dependence the observable has. For example, a GRACE coefficient averaged over one
month has $\zeta^\dagger_\phi(t)$ equal to the spatial pattern of the third table row multiplied by
a boxcar in time over that month; a rate estimated by regression over several years has
$\zeta^\dagger(t)$ proportional to the regression weights. The adjoint problem is then solved
backwards in time as in eq.~\eqref{d319a}.

\subsection{The equatorial equation in a spherically symmetric model}
\label{sec:dyn_spec}

With $I_1=I_2$, the equatorial part of $\bG$ is $(I_1/\Omega)\E$, where
\begin{equation*}
  \E=\begin{pmatrix}0&-1\\1&0\end{pmatrix}
\end{equation*}
rotates a vector by $+90^\circ$ and satisfies $\E^T=-\E$ and $\E^2=-\Id$. With $\bq$ defined as in
eq.~\eqref{p3}, the equatorial part of eq.~\eqref{d320} reads
$(I_3-I_1)\bw_\perp+(I_1/\Omega)\E\dot\bw_\perp=(I_3-I_1)\bq$, that is
\begin{equation}
  \boxed{\;\dot\bw_\perp=\sigma_e\,\E\,(\bw_\perp-\bq),\qquad \sigma_e=\frac{\Omega(I_3-I_1)}{I_1},
  \qquad \bq=\beta\bm\phi_{21}-\frac{\bk_\perp}{I_3-I_1}\;}
  \tag{DP1}\label{dp1}
\end{equation}
or, in components,
\begin{equation}
  \dot\omega_1=-\sigma_e(\omega_2-q_2),\qquad \dot\omega_2=+\sigma_e(\omega_1-q_1).
  \tag{DP1a}\label{dp1a}
\end{equation}

\begin{clar}
To get eq.~\eqref{dp1}, divide the equatorial equation by $I_1/\Omega$, which gives
$\E\dot\bw_\perp=\sigma_e(\bq-\bw_\perp)$, then apply $\E^{-1}=-\E$. The static relation
eq.~\eqref{p3} is recovered by setting $\dot\bw_\perp=\mathbf{0}$, since $\E$ is invertible.
\end{clar}

\paragraph{Reading eq.~\eqref{dp1}.} The static theory says $\bw_\perp=\bq$: the rotation vector
sits at the forcing. The dynamic theory says that $\bw_\perp-\bq$ is rotated by $\E$ and fed back
into $\dot\bw_\perp$, which means the rotation vector \emph{circles around} $\bq$ rather than
sitting on it. For constant $\bq$ the solution is exactly
\begin{equation}
  \bw_\perp(t)=\bq+\R(\sigma_e t)\,[\bw_\perp(0)-\bq],
  \tag{DP2}\label{dp2}
\end{equation}
with $\R(\alpha)$ the rotation by angle $\alpha$. So the motion is a circle of radius
$|\bw_\perp(0)-\bq|$ traced at rate $\sigma_e$, in the same sense as the Earth's rotation
(prograde). The static solution is the \emph{centre} of that circle. It is recovered when the radius
is negligible, which happens when $\bq$ changes slowly compared with the circling: eq.~\eqref{dp7}
below shows that a steadily moving $\bq$ keeps the radius at roughly $|\dot\bq|/\sigma_{\rm CW}$,
i.e.\ the distance $\bq$ travels in one radian of wobble.

\paragraph{The feedback, and the true wobble frequency.} $\bq$ is not an external quantity: part of
it is generated by $\bw_\perp$ itself, through the deformation raised by $\psi$ and the equilibrium
ocean pole tide. By linearity,
\begin{equation}
  \bq(t)=\bq^{\zeta}(t)+\K\,\bw_\perp(t),
  \tag{DP3}\label{dp3}
\end{equation}
where $\bq^\zeta$ is driven by the direct load (with the ocean responding) and by $\bk$, and $\K$ is
a constant $2\times2$ matrix; without oceans $\K=\kappa\Id$ with $\kappa=0.314$ from
eq.~\eqref{p4}. Substituting,
\begin{equation}
  \dot\bw_\perp=\sigma_e\E\big[(\Id-\K)\bw_\perp-\bq^\zeta\big],
  \tag{DP4}\label{dp4}
\end{equation}
so the free motion circles at
\begin{equation}
  \sigma_{\rm CW}=\sigma_e\sqrt{\det(\Id-\K)}\;\xrightarrow[\text{no ocean}]{}\;\sigma_e(1-\kappa),
  \qquad T_{\rm CW}=\frac{304.5\,\mathrm{d}}{1-0.314}\approx444\,\mathrm{d},
  \tag{DP5}\label{dp5}
\end{equation}
against an observed 433\,d. Deformation lengthens the period because the Earth's bulge partly
follows the displaced rotation axis, weakening the restoring effect.

\begin{clar}
For symmetric $\K$ the matrix $\sigma_e\E(\Id-\K)$ has zero trace, so its eigenvalues are
$\pm i\sigma_{\rm CW}$: the free motion is undamped, and slightly elliptical if $\K$ is anisotropic
(which the ocean geometry makes it). A non-zero trace in a numerical $\K$ would produce spurious
growth or decay, and is a useful check. Symmetry of $\K$ is expected on reciprocity grounds but I
have not proved it.
\end{clar}

\subsection{How the dynamic response differs from the static one}
\label{sec:dyn_resp}

\paragraph{Periodic loading.} Write the equatorial vector as a complex number,
$w=\omega_1+i\omega_2$, so that $\E$ becomes multiplication by $i$. A circular motion
$w\propto e^{i\nu t}$ with $\nu>0$ turns in the same sense as the Earth (prograde) and with $\nu<0$
in the opposite sense (retrograde); any real forcing is a sum of the two. For forcing at frequency
$\nu$ and $\K=\kappa\Id$, eq.~\eqref{dp4} gives
\begin{equation}
  \frac{w_{\rm dynamic}}{w_{\rm static}}=\frac{1}{1-\nu/\sigma_{\rm CW}} ,
  \tag{DP6}\label{dp6}
\end{equation}
where $w_{\rm static}$ is what AA24's theory would give for the same load. This is a resonance at
the Chandler frequency, and only in the prograde sense:
\begin{center}
\begin{tabular}{lcccc}
\toprule
Forcing period & 30\,d & 182.6\,d & 365.25\,d & $\to\infty$\\
\midrule
prograde & $-0.072$ & $-0.70$ & $-4.64$ & $1$\\
retrograde & $0.063$ & $0.29$ & $0.45$ & $1$\\
\bottomrule
\end{tabular}
\end{center}
Reading the table: at monthly periods the polar motion produced by a given load is an order of
magnitude \emph{smaller} than quasi-static theory predicts, and reversed in sign, because the
rotation axis has no time to respond. At annual periods the prograde part is amplified about
fivefold and again reversed. Only for forcing slow compared with 444\,d does the static answer
apply. Since the rotational feedback is itself a small part of the sea-level response, these
factors multiply a small correction; they say how wrong the correction is, not how wrong sea level
is.

\begin{clar}
The numbers use $T_{\rm CW}=444$\,d and $\K=\kappa\Id$, so they omit the ocean, core and anelastic
corrections to $\sigma_{\rm CW}$. The annual prograde entry sits close to the resonance and is
therefore sensitive to that value; it should not be taken literally.
\end{clar}

\paragraph{Secular loading.} For a load whose $\bq^\zeta$ grows linearly from zero,
$\bq^\zeta=\dot\bq^\zeta t$, the exact solution of eq.~\eqref{dp4} is
\begin{equation}
  \bw_\perp(t)=\underbrace{\frac{\dot\bq^\zeta}{1-\kappa}\,t}_{\text{quasi-static}}
  +\underbrace{\frac{1}{\sigma_{\rm CW}}\E^{T}\frac{\dot\bq^\zeta}{1-\kappa}}_{\text{constant offset}}
  +\underbrace{\bw_{\rm free}(t)}_{\text{Chandler wobble}},
  \tag{DP7}\label{dp7}
\end{equation}
where $\bw_{\rm free}$ is a free circling motion whose amplitude equals that of the offset term,
fixed by $\bw_\perp(0)=\mathbf{0}$. The trend itself is exactly the quasi-static one. The offset is
perpendicular to the trend and equals the trend rate multiplied by
$T_{\rm CW}/2\pi\approx70$\,d, i.e.\ the polar motion lags the static prediction by about 70 days.
For multi-year trends this is a small correction to an accumulated signal, so trend-based
inferences, such as those in Heathcote et al., are unlikely to be affected. Monthly anomalies are
the case where it matters.

\subsection{Time stepping and iterative solution}
\label{sec:dyn_num}

The practical consequence of all the above is modest: the solver gains one extra state variable,
the 2-vector $\bw_\perp$, which is carried from one time step to the next instead of being
recomputed from scratch.

\paragraph{The step map.} Applying the trapezoidal rule to eq.~\eqref{dp1} over a step $\Delta t$,
with $\bw_n=\bw_\perp(t_n)$, $\bq_n=\bq(t_n)$ and $a=\sigma_e\Delta t/2$, gives
\begin{equation}
  \boxed{\;\bw_{n+1}=\R\,\bw_n+\Sm\,(\bq_n+\bq_{n+1}),\qquad
  \R=\frac{(1-a^2)\Id+2a\E}{1+a^2},\qquad \Sm=\frac{a^2\Id-a\E}{1+a^2}\;}
  \tag{DP8}\label{dp8}
\end{equation}
$\R$ is exactly a rotation by $2\arctan a$: it carries the free wobble forward without changing its
amplitude, so the scheme adds no artificial damping. $\Sm$ is the weight with which the forcing
pulls $\bw$ towards the static solution over the step. Two limits are worth checking:
if $\bq$ is constant and $\bw_n=\bq$, then $\bw_{n+1}=\bq$, since $\R+2\Sm=\Id$; and for
$a\to\infty$ the rotation angle tends to $\pi$, which is the (non-L-stable) way the trapezoidal rule
represents an infinitely fast wobble. If instead a scheme that collapses onto the static theory for
large steps is wanted, backward Euler with $b=\sigma_e\Delta t$ gives
\begin{equation}
  \bw_{n+1}=\frac{(\Id+b\E)\,\bw_n+(b^2\Id-b\E)\,\bq_{n+1}}{1+b^2}\;\xrightarrow[b\to\infty]{}\;\bq_{n+1},
  \tag{DP9}\label{dp9}
\end{equation}
at the cost of damping the wobble numerically. For monthly steps, $a\approx0.31$ and the trapezoidal
free period is too long by 1.5\% (450.6\,d against 444.0\,d).

\paragraph{Scheme A: modify the existing iteration.} The static solver already iterates, at each
time, between the load, the fields and $\bw$. Only the $\bw$ update changes. Given $\bw_n$, $\bq_n$
and the load $\zeta_{n+1}$:
\begin{center}
\fbox{\begin{minipage}{0.93\linewidth}
\begin{enumerate}
  \item Initialise $\bw^{(0)}=\bw_n$ and $\sigma^{(0)}=\zeta_{n+1}+\rho_w\oc\,\dSL_n$.
  \item For $k=0,1,\dots$ until converged:
  \begin{enumerate}
    \item $\bm\psi_{21}=r\,\bw^{(k)}$ from eq.~\eqref{p1};
    \item $u_{lm}=h_l\sigma^{(k)}_{lm}$ and $\phi_{lm}=k_l\sigma^{(k)}_{lm}$, adding
    $h^T_2\bm\psi_{21}$ and $k^T_2\bm\psi_{21}$ at $(l,m)=(2,1)$;
    \item $\bq^{(k)}=\beta\bm\phi_{21}-\bk_{\perp,n+1}/(I_3-I_1)$;
    \item \textbf{the only change:} $\bw^{(k+1)}=\R\bw_n+\Sm(\bq_n+\bq^{(k)})$, in place of the
    static update $\bw^{(k+1)}=\bq^{(k)}$;
    \item add $\psi$ from $\bw^{(k+1)}$ to $\phi$, form $\dSL$ from eq.~\eqref{aa318}, and update
    $\sigma^{(k+1)}=\zeta_{n+1}+\rho_w\oc\,\dSL$.
  \end{enumerate}
  \item Store $\bw_{n+1}$ and $\bq_{n+1}$ (the latter with $\phi$ excluding $\psi$) for the next step.
\end{enumerate}
\end{minipage}}
\end{center}
The static solver is the special case $\R=0$, $\Sm=\Id$, with the $\bq_n$ term dropped. Convergence
is no worse than in the static case: the rotational part of the iteration contracts by at most
$\|\Sm\|\,\|\K\|=\kappa a/\sqrt{1+a^2}<\kappa$.

\paragraph{Scheme B: solve the rotation separately.} Because the elastic Earth and the equilibrium
ocean respond instantaneously, $\bw$ can be eliminated exactly, which avoids repeating the
sea-level solve every time the wobble is stepped:
\begin{enumerate}
  \item \emph{Once:} for $\alpha=1,2$, solve the sea-level equation with no direct load and a
  \emph{prescribed} $\bw_\perp=\mathbf{e}_\alpha$ (impose $\psi$, let the ocean respond, but do not
  update $\bw$). This gives the response fields $X^\alpha$ and the columns of $\K$.
  \item \emph{For each time:} solve the sea-level equation for $\zeta(t)$ with $\bw_\perp\equiv\mathbf{0}$,
  i.e.\ with rotational feedback switched off. This gives $X^\zeta(t)$ and $\bq^\zeta(t)$.
  \item \emph{Then:} integrate the $2\times2$ problem eq.~\eqref{dp4} on its own, either exactly or
  with eq.~\eqref{dp8}, whose implicit step needs only the $2\times2$ solve
  $(\Id-\Sm\K)\bw_{n+1}=\R\bw_n+\Sm(\bq_n+\bq^\zeta_{n+1})$, where $\bq_n=\bq^\zeta_n+\K\bw_n$ is
  already known.
  \item \emph{Finally:} reassemble the full solution as
  $X(t)=X^\zeta(t)+\sum_\alpha\omega_\alpha(t)\,X^\alpha$.
\end{enumerate}
This costs two extra sea-level solves in total, allows the wobble to be sub-stepped as finely as
desired, and produces $\K$ explicitly for the checks in \S\ref{sec:dyn_checks}. Scheme~A is the
smaller change to existing code; Scheme~B is better for long time series and for inversions.

\paragraph{Initial conditions.} If the model starts from rest, $\bw_0=\mathbf{0}$. If it starts in
static equilibrium under a load already present, set $\bw_0$ to the static solution. A free wobble
present at the start is not determined by the loads at all; it is either prescribed or, as in
eq.~\eqref{d319b}, estimated.

\paragraph{Adjoint.} By eq.~\eqref{d319a}, adjoint fields are computed with the same solver applied
to time-reversed adjoint loads and $\sigma_e\to-\sigma_e$, i.e.\ $a\to-a$, which simply replaces
$\R$ and $\Sm$ by $\R^T$ and $\Sm^T$. That is the continuous adjoint; for dot-product tests at
machine precision, the discrete adjoint should instead be obtained by transposing the recursion
eq.~\eqref{dp8}, which agrees with it to $O(\Delta t^2)$.

\subsection{Summary of the recipe}

Starting from a working static solver, the changes are:
\begin{enumerate}
  \item Carry $\bw_\perp$ as a state variable, one 2-vector per time step, instead of recomputing it
  from scratch at each time.
  \item At each time step, run the existing iteration unchanged except for the rotation update,
  which becomes eq.~\eqref{dp8} in place of $\bw_\perp\leftarrow\bq$. The matrices $\R$ and $\Sm$
  depend only on $a=\sigma_e\Delta t/2$, and $\sigma_e=\Omega(I_3-I_1)/I_1$ is a constant of the
  earth model.
  \item Or, equivalently and more cheaply for long runs, switch the rotational feedback off
  entirely, solve the $2\times2$ wobble problem separately, and superpose (Scheme~B).
  \item For adjoints, reverse the loads in time and replace $\R,\Sm$ by $\R^T,\Sm^T$.
\end{enumerate}
Nothing else changes: Love numbers, the ocean iteration, mass conservation and the form of $\psi$
are all as in the static theory. The physical content is eq.~\eqref{dp2} --- the rotation vector
circles around the static answer at the Chandler rate --- and the practical consequence is the table
in \S\ref{sec:dyn_resp}: the static theory is right for trends, and wrong, by a factor of order ten
at monthly periods, for sub-annual variability.

\subsection{Checks and open issues}
\label{sec:dyn_checks}

\paragraph{Checks worth running.}
\begin{enumerate}
  \item Backward Euler with a very large $\Delta t$ reproduces the existing static \texttt{pyslfp}
  solution.
  \item With no load and no ocean, and $\bw_0\neq\mathbf{0}$, the solution circles with period
  444\,d at constant $|\bw_\perp|$.
  \item $\K$ from Scheme~B is symmetric, with zero trace in $\E\K$, and reduces to $\kappa\Id$ when
  the ocean function is set to zero.
  \item Schemes A and B agree for the same $\Delta t$.
  \item A space--time dot-product test of eq.~\eqref{d319}, forward against time-reversed adjoint.
\end{enumerate}

\paragraph{Physics not included.}
\begin{itemize}
  \item \textbf{Fluid core.} It is largely decoupled from the wobble, so the inertial term should
  involve mantle moments and wobble-appropriate Love numbers (Smith \& Dahlen 1981). Pragmatically,
  $\sigma_e$ can be tuned so that $T_{\rm CW}$ matches the observed 433\,d.
  \item \textbf{Anelasticity.} It gives $Q_{\rm CW}\sim10^2$. A damping term breaks the
  time-reversal structure, so the adjoint then carries damping of the opposite sign.
  \item \textbf{Triaxiality} ($I_1\neq I_2$). Eq.~\eqref{d37a} fails, so the criterion of
  \S\ref{sec:dyn_jdot} --- every matrix multiplying $\dd/\dd t$ must be a multiple of $\E$ --- is
  violated and reciprocity in the form of eq.~\eqref{d319} is not established here. The same
  criterion governs the ocean term in eq.~\eqref{d312}.
  \item \textbf{Observational conventions.} GRACE and altimetry pole-tide corrections are made using
  observed polar motion, and the model treatment should mirror that processing (cf.\ Lickley et al.\
  2018). This is the dynamic version of the question raised in \S\ref{sec:gen_rot} about which
  centrifugal terms belong in the observable.
\end{itemize}

\begin{thebibliography}{9}
\bibitem{aa24} Al-Attar, D., Syvret, F., Crawford, O., Mitrovica, J.X.\ \& Lloyd, A.J., 2024. Reciprocity and sensitivity kernels for sea level fingerprints, \emph{Geophys.\ J.\ Int.}, 236(1), 362--378. arXiv:2307.10835.
\bibitem{at14} Al-Attar, D.\ \& Tromp, J., 2014. Sensitivity kernels for viscoelastic loading based on adjoint methods, \emph{Geophys.\ J.\ Int.}, 196(1), 34--77.
\bibitem{dahlen74} Dahlen, F.A., 1974. On the static deformation of an earth model with a fluid core, \emph{Geophys.\ J.\ R.\ astr.\ Soc.}, 36, 461--485.
\bibitem{lickley} Lickley, M.J., Hay, C.C., Tamisiea, M.E.\ \& Mitrovica, J.X., 2018. Bias in estimates of global mean sea level change inferred from satellite altimetry, \emph{J.\ Climate}, 31(13), 5263--5271.
\bibitem{mitrovica05} Mitrovica, J.X., Wahr, J., Matsuyama, I.\ \& Paulson, A., 2005. The rotational stability of an ice-age earth, \emph{Geophys.\ J.\ Int.}, 161(2), 491--506.
\bibitem{munk} Munk, W.H.\ \& MacDonald, G.J.F., 1960. \emph{The Rotation of the Earth}, Cambridge University Press.
\bibitem{sd81} Smith, M.L.\ \& Dahlen, F.A., 1981. The period and Q of the Chandler wobble, \emph{Geophys.\ J.\ R.\ astr.\ Soc.}, 64, 223--281.
\bibitem{tm99} Tromp, J.\ \& Mitrovica, J.X., 1999. Surface loading of a viscoelastic earth---I. General theory, \emph{Geophys.\ J.\ Int.}, 137, 847--855.
\end{thebibliography}

\end{document}
